% ============================================================
\section{Experiment 2: Does Encoder Adaptation Help or Harm?}
\label{sec:exp2}
% ============================================================

\subsection{What Fine-Tuning Does to the Seven Survivors}
\label{sec:forecasting:results}

The seven cells Experiment 1 admits are the whole evidence base for a claim about
damage to valuable features, and they do not support one
(Table~\ref{tab:gate_survivors}).

\textbf{On four of the seven, full fine-tuning improves the checkpoint in every
seed.}  On Moirai-Base at $h{=}96$ and $h{=}192$ and Moirai-Large at $h{=}96$, all
on ETTh2, forgetting is $-31.4\%$, $-36.3\%$ and $-42.4\%$ against zero-shot,
negative in 3/3 seeds each, at CKA $0.460$, $0.396$ and $0.626$, the largest
representational change among the seven.  Freezing the encoder captures only part
of that gain ($-9.1\%$, $-16.1\%$, $-21.6\%$), so most of the improvement comes
from encoder adaptation.  The fourth, Moirai-Base/ETTm2 at $h{=}192$, improves by
$8.4\%$, and the frozen encoder does slightly better there ($-9.7\%$).

\textbf{On the remaining three, the mean harm does not survive inspection.}
Moirai-Small/ETTh2 at $h{=}192$ reads forg.$_\text{B}{=}{+}1.99{\pm}3.32$ with 5 of
10 seeds positive, the closest the matrix comes to degradation.  At $h{=}96$ the
positive mean ($+9.71{\pm}10.58$, 4 of 10 seeds) comes from one seed that diverges
under the frozen encoder too, an optimisation failure rather than an encoder
effect; we keep it in every mean.  Moirai-Small/ETTm2 at $h{=}192$ improves on
average ($1.10\%$) with 3 of 5 seeds harmful.

\subsection{Where the Freezing Boundary Is Drawn}
\label{sec:exp2:strict}

The ordinary frozen encoder (D) leaves the input projection and mask encoding
trainable, so a ``freezing wins'' reading could come from those modules adapting
rather than from the encoder staying fixed.  We therefore ran the strict freeze
(H) on every Moirai cell (\pending{22} cells, 3--10 seeds;
Table~\ref{tab:strictfreeze}).  B${-}$H and B${-}$D agree in sign on
\pending{N} cells.  They disagree on \pending{the four Moirai/ETTm2 cells
and any of the six new ones}, and on every disagreeing cell the strict control
moves toward adaptation: on ETTm2, Moirai-Small $h{=}96$ goes from B${-}$D
$+7.3{\pm}2.9$ to B${-}$H $-6.8{\pm}2.0$ and Moirai-Base $h{=}96$ from $+7.0{\pm}2.4$
to $-9.2{\pm}3.1$, so the apparent freezing advantage on that dataset was carried
by the input projection.  The strict freeze therefore recovers no degradation case;
it removes the ones D appeared to suggest.  All Moirai counts below are reported
under both controls.

\subsection{Held-Out Scoring Is a Precondition}
\label{sec:analysis:heldout}

Scoring $\denc$ on the selection windows instead of the held-out ones
\textbf{reverses its sign in 4 of the 31 intervention cells}, all in the same
direction: on the selection windows the frozen encoder wins, on held-out windows
full fine-tuning does (Table~\ref{tab:heldout_decomp}).  Two of the four are
among the seven gate-passing cells (Moirai-Base and Moirai-Large on ETTh2, which
move from $+3.8$ and $+7.2$ to $-22.3$ and $-20.8$).  A frozen-versus-full comparison
scored on the windows a checkpoint was selected on is therefore not a weaker
version of the held-out one; it can give the opposite answer.

\subsection{No Degradation Instance Survives a Prespecified Definition}
\label{sec:analysis:degradation}

The three-clause definition of \S\ref{sec:method} returns \textbf{zero cells} of
the 31, under both the interval and the registered form and, on Moirai, under
both controls.  Clause~(i) removes 24 outright: they have no measured pre-trained
advantage, so a worse score after fine-tuning replaces nothing demonstrated.
Of the seven that remain, four fail clause~(ii) with every seed agreeing, and
the other three fail it cell by cell (\S\ref{sec:forecasting:results}).
Relaxing the definition does not manufacture a case: across 8 rungs, at either
gate threshold, \textbf{the largest count anywhere is 2}, reached only when the
means clause is dropped altogether (Appendix~\ref{app:defsensitivity}).

\textbf{Per-cell inference.}  Paired $95\%$ intervals on $\denc$ (F${=}$D),
Benjamini--Hochberg adjusted across the 31 cells, call \textbf{2 cells
freezing-better, 4 adaptation-better and 25 inconclusive}; under F${=}$H on
Moirai, \pending{counts}.  The two freezing-better cells are Moirai-Small/ETTh1 at
both horizons, which the screen rejects, while two of the four adaptation-better
cells pass it, at $+0.925$ and $+0.875$.  A pre-registered seed top-up adds a third
freezing-better cell, which the screen admits (Moirai-Small/ETTm2 $h{=}192$,
gate $+0.352$), and the definition still returns zero
(Appendix~\ref{app:pairedinference}).  Most calls are inconclusive because the
design is underpowered for small effects, and we say so rather than read them as
null: the survivors' dispersion gives minimum detectable effects of
$4.7$--$14.7$~pp, and two survivors would need 182 and 212 paired seeds to
resolve their observed effect.  Over the 20 value-cells, $\denc{=}{+}0.7$ with
clustered CI $[-11.1,{+}14.1]$.  Changing the gate's estimator does admit cells
(persistence admits three and GBM one), but all four rest on a denominator within
$5\%$ of the training-mean floor, so a weaker baseline is what manufactures them
(Appendix~\ref{app:baselines}).

We do not conclude that fine-tuning never damages a pre-trained capability, only
that on these benchmarks we could not find one instance that meets its own
definition: no cell's advantage survives every admissible rung of the ladder, and
where the fitted-ridge advantage exists, fine-tuning on 500--1{,}000 samples
usually improves on it.
